\chapter{Struttura delle directory}

La struttura principale del repository è la seguente:

\begin{lstlisting}
PlantPulse/
|-- .github/workflows/ci.yml
|-- dashboard/
|   |-- backend/
|   |   |-- app/
|   |   |   |-- main.py
|   |   |   |-- database.py
|   |   |   |-- models.py
|   |   |   |-- schemas.py
|   |   |   `-- logging_config.py
|   |   |-- sql/init.sql
|   |   |-- tests/
|   |   `-- Dockerfile
|   |-- frontend/
|   |   |-- src/components/
|   |   |-- src/pages/
|   |   |-- src/services/api.js
|   |   |-- Dockerfile
|   |   `-- nginx.conf
|   |-- docker-compose.yml
|   `-- docker-compose.prod.yml
|-- hardware/main.cpp
|-- .gitignore
`-- README.md
\end{lstlisting}

\section{Responsabilità dei file principali}
\begin{itemize}
    \item \texttt{main.py}: definizione delle route FastAPI e logica applicativa degli endpoint.
    \item \texttt{database.py}: caricamento della configurazione, engine SQLAlchemy e sessioni.
    \item \texttt{models.py}: modelli ORM delle tre tabelle.
    \item \texttt{schemas.py}: validazione e serializzazione dei dati con Pydantic.
    \item \texttt{init.sql}: schema iniziale PostgreSQL, vincoli, indice e configurazione iniziale del dispositivo.
    \item \texttt{hardware/main.cpp}: simulazione, invio misurazioni, recupero comandi e acknowledgement.
    \item \texttt{frontend/src/services/api.js}: funzioni client per le chiamate HTTP.
    \item \texttt{ci.yml}: pipeline di verifica e pubblicazione delle immagini.
\end{itemize}
